% ===============================
% Worksheet / Manual Template
% ===============================
\documentclass[12pt, a4paper]{article}

% ===============================
% Metadata
% ===============================
\newcommand*{\CourseCode}{MAT321}
\newcommand*{\CourseTitle}{Real Analysis II}
\newcommand*{\Chapter}{Subspaces}
\newcommand*{\Topics}{Interior, Exterior, Boundary and Subspaces}
\newcommand*{\DocDate}{February 17, 2026}

\include{header}

% ==========================================
% Document
% ==========================================
\begin{document}
\FrontPage
\Large

\begin{heading}
    Interior, Exterior, Frontier and Boundary Points
\end{heading}

\vspace{10pt}

\begin{definition}[Interior Point]
Let $A$ be a subset of a metric space $(X,d)$. A point $a\in A$ is called an
\emph{interior point} of $A$ if there exists $r>0$ such that
\[
a\in S_r(a)\subseteq A.
\]
\end{definition}

\vspace{150pt}

\begin{definition}[Interior of a Set]
The set of all interior points of $A$ is called the \emph{interior} of $A$ and is denoted by
$\operatorname{int}A$, where
\[
\operatorname{int}A=\{x\in A : S_r(x)\subseteq A \text{ for some } r>0\}.
\]
Clearly, $\operatorname{int}A$ is an open subset of $A$.
\end{definition}

\clearpage

\begin{definition}[Exterior Point]
A point $x\in X$ is called an \emph{exterior point} of $A$ if it is an interior point of the
complement of $A$, that is, if there exists $r>0$ such that
\[
S_r(x)\subseteq A^c
\quad\text{or equivalently}\quad
S_r(x)\cap A=\emptyset.
\]
\end{definition}

\vspace{150pt}

\begin{definition}[Exterior of a Set]
The set of all exterior points of $A$, denoted by $\operatorname{ext}A$, is called the
\emph{exterior} of $A$. By definition,
\[
\operatorname{ext}A=\operatorname{int}A^c.
\]
Hence $\operatorname{ext}A$ is open and is the largest open set contained in $A^c$.
Moreover,
\[
\operatorname{ext}A=(\overline{A})^c.
\]
\end{definition}

\clearpage

\begin{definition}[Boundary Point]
A point $x\in X$ is called a \emph{boundary point} of $A$ if
every open sphere centered at $x$ intersects both $A$ and its complement $A^c$, that is,
for every $r>0$,
\[
S_r(x)\cap A\neq\emptyset
\quad\text{and}\quad
S_r(x)\cap A^c\neq\emptyset.
\]
\end{definition}

\vspace{150pt}

\begin{definition}[Boundary of a Set]
The set of all boundary points of $A$ is called the \emph{boundary} of $A$
and is denoted by $\partial A$, where
\[
\partial A=\overline{A}\setminus \operatorname{int}A.
\]
Equivalently,
\[
\partial A=\overline{A}\cap \overline{A^c}.
\]
\end{definition}

\clearpage

\begin{problem}[Illustrations of Interior, Exterior, Boundary]
Compute the interior, exterior, and boundary of the following subsets.

\begin{enumerate}
    \item Let $X=\mathbb{R}$ with the usual metric and let $A=\mathbb{Q}$. Then
    \[
    \operatorname{int}A=\emptyset,\qquad
    \operatorname{ext}A=\emptyset,\qquad
    \partial A=\mathbb{R}.
    \]

    \item Let $X=\mathbb{R}$ with the usual metric and let
    \(
    A=(1,2)\cup(3,4).
    \)
    Then
    \[
    \operatorname{int}A=A,
    \]
    \[
    \operatorname{ext}A=(-\infty,1)\cup(2,3)\cup(4,\infty),
    \]
    \[
    \partial A=\{1,2,3,4\}.
    \]

    \item If $(X,d)$ is a discrete metric space and $A\subseteq X$, then
    \[
    \operatorname{int}A=A,\qquad
    \operatorname{ext}A=A^c,\qquad
    \partial A=\emptyset.
    \]

    \item Let $X=\mathbb{N}$ with the usual metric and let $A=\{1,2,3\}$. Then
    \[
    \operatorname{int}A=\emptyset,\qquad
    \operatorname{ext}A=\emptyset,\qquad
    \partial A=\mathbb{N}.
    \]
\end{enumerate}
\end{problem}

\clearpage

\begin{theorem}[Properties of Interior]
Let $A$ and $B$ be any two subsets of a metric space $(X,d)$. Then:
\begin{enumerate}
    \item $\operatorname{int}A$ is the largest open set contained in $A$, that is,
    \[
    \operatorname{int}A=\bigcup\{G : G \text{ is open and } G\subseteq A\};
    \]
    
    \item $A$ is open if and only if $A=\operatorname{int}A$;
    
    \item $A\subseteq B$ implies $\operatorname{int}A\subseteq \operatorname{int}B$;
    
    \item $\operatorname{int}(A\cap B)=\operatorname{int}A\cap \operatorname{int}B$;
    
    \item $\operatorname{int}(A\cup B)\supseteq \operatorname{int}A\cup \operatorname{int}B$.
\end{enumerate}
\end{theorem}


\clearpage

\begin{theorem}[Properties of Exterior of a Set]
Let $(X,d)$ be a metric space and let $A,B$ be any two subsets of $X$. Then:
\begin{enumerate}
    \item $\operatorname{ext}A$ is the largest open set contained in $A^c$;
    \item $A$ is closed if and only if $A^c=\operatorname{ext}A$;
    \item $A\subseteq B$ implies $\operatorname{ext}B\subseteq \operatorname{ext}A$;
    \item $\operatorname{ext}(A\cap B)\supseteq \operatorname{ext}A\cup \operatorname{ext}B$;
    \item $\operatorname{ext}(A\cup B)=\operatorname{ext}A\cap \operatorname{ext}B$.
\end{enumerate}
\end{theorem}

\clearpage

\begin{theorem}[Properties of Boundary]
Let $(X,d)$ be a metric space, and let $A,B$ be subsets of $X$. Then:
\begin{enumerate}
    \item $\partial A=\overline{A}\cap \overline{A^c}
    =\overline{A}\setminus \operatorname{int}A$;
    
    \item $\partial A=\emptyset$ if and only if $A$ is both open and closed;
    
    \item $A$ is closed if and only if $\partial A \subseteq A$;
    
    \item $A$ is open if and only if $\partial A \subseteq A^c$;
    
    \item $\partial (A\cap B)\subseteq \partial A\cup \partial B$.
    The equality holds if $\overline{A}\cap \overline{B}=\emptyset$;
    
    \item $\partial(\operatorname{int}A)\subseteq \partial A$.
\end{enumerate}
\end{theorem}

\clearpage


\begin{heading}
    Subspaces
\end{heading}

\vspace{10pt}

\begin{definition}[Subspace and Induced Metric]
Let $(X,d)$ be a metric space and let $Y$ be a non-empty subset of $X$.
The restriction of the metric $d$ to $Y\times Y$, denoted by $d_Y$, is a metric on $Y$
called the \emph{induced metric}. The metric space $(Y,d_Y)$ is called a
\emph{subspace} of $(X,d)$.
\end{definition}

\vspace{100pt}

\begin{problem}[Examples of Subspaces]
Verify that:
\begin{enumerate}
    \item the closed unit interval $[0,1]$ and the set of all rational numbers $\mathbb{Q}$
    are subspaces of $\mathbb{R}$ with the usual metric;
    \item the unit circle, the closed unit disc, and the open unit disc are subspaces of
    the metric space $(\mathbb{C},d)$ of complex numbers;
    \item the real line $\mathbb{R}$ itself is a subspace of the complex plane $\mathbb{C}$.
\end{enumerate}
\end{problem}

\clearpage

\begin{definition}[Open Sphere in a Subspace]
Let $(X,d)$ be a metric space, let $Y\subseteq X$, and let $y\in Y$.
The open sphere in the subspace $(Y,d_Y)$ centered at $y$ with radius $r>0$ is denoted by
$S_r^{\,Y}(y)$ and is defined by
\[
S_r^{\,Y}(y)=\{x\in Y : d_Y(x,y)<r\}.
\]
\end{definition}

\vspace{100pt}

\begin{problem}[Relation Between Open Spheres]
Show that for every $y\in Y$ and every $r>0$,
\[
S_r^{\,Y}(y)=S_r(y)\cap Y,
\]
where $S_r(y)$ denotes the open sphere in $(X,d)$.
\end{problem}


\clearpage

\begin{theorem}[Open Sets in Subspaces]
Let $(X,d)$ be a metric space and let $Y\subseteq X$. A subset $A$ of $Y$ is open in
$(Y,d_Y)$ if and only if there exists a set $G$ open in $(X,d)$ such that
\[
A = G \cap Y.
\]
\end{theorem}

\clearpage

\begin{theorem}[Open Sets in a Subspace vs Ambient Space]
Show that if $Y\subseteq X$ and a subset of $Y$ is open in $(X,d)$, then it is also open in
the subspace $(Y,d_Y)$. Show by examples that the converse need not be true.
\end{theorem}

\vspace{20pt}

\begin{problem}[Open in a Subspace may Not be open in Superspace]
Let $X=\mathbb{R}$ with the usual metric and let
\[
Y=[0,1].
\]
Then the open sphere in $(Y,d_Y)$ centered at $0$ with radius $\tfrac12$ is
\[
S^{\,Y}_{1/2}(0)=[0,\tfrac12),
\]
which is open in $Y$ but not open in $\mathbb{R}$. Note that $Y$ itself is not open in
$\mathbb{R}$.
\end{problem}

\vspace{80pt}

\begin{problem}[Example: Real Line vs Complex Plane]
Show that an open interval of the real line is not an open subset of the complex plane
$\mathbb{C}$ with the usual metric.
\end{problem}

\clearpage

\EndPage
\end{document}